\section{Conclusiones}
\label{sec:conclusiones}

\subsection*{Hallazgos cuantitativos}

Sobre $5\,750$ clientes, $20\,000$ tickets y $42\,555$ líneas de venta, el
pipeline ETL puebla un DWH en estrella ($6$~dimensiones, $2$~hechos) de
forma reproducible: una sola invocación de
\texttt{jupyter nbconvert --execute notebooks/0*} regenera los esquemas
\texttt{staging} y \texttt{dwh}, los \emph{parquet} de analítica y los
modelos de \texttt{models/}. La distribución del CLTV es extremadamente
asimétrica (mediana $0{,}72$\,\texteuro\ frente a media $3\,165$\,\texteuro,
máximo $109\,735$\,\texteuro), con un coeficiente de Gini~$\approx 0{,}91$:
el $87\%$ de los clientes son \emph{one-shot} y el cluster ``alto valor''
($13\%$ de la base) concentra el grueso del CLTV agregado. El segmentador
PCA(2) $+$ KMeans(3) recupera tres clusters interpretables (ocasionales,
alto valor, devolutivos) con $80{,}3\%$ de varianza explicada y un codo
claro en $k^{\star}=3$.

\subsection*{Limitaciones}

El dataset es \emph{sintético} y algunas decisiones de diseño preservan los
datos tal cual aunque generen artefactos: existen $10$ clientes con
$\texttt{return\_rate} > 1$ (devoluciones que superan la venta original) y
$1$ producto sin \texttt{unit\_cost} en \texttt{central\_product}; ambos
casos se conservan y se gestionan con \texttt{coalesce(\dots, 0)} en el
cálculo de coste. La retención $R_t$ es una proxy discreta
($\texttt{active\_months}/72$) y el KMeans se entrena sobre el espacio
reducido por PCA(2), que descarta un $\sim 20\%$ de la varianza original.

\subsection*{Líneas futuras y validación}

Las extensiones naturales son: ampliar las \emph{features} con señales
territoriales (\texttt{district}, \texttt{area\_type}) y promocionales
(\texttt{discount\_percent}); validar $k$ con métricas de silueta y
Davies--Bouldin; sustituir la fórmula multiplicativa por modelos
probabilísticos tipo BG/NBD + Gamma--Gamma; y refactorizar el pipeline
para ingestión incremental con \emph{slowly changing dimensions}. La
calidad queda anclada por nueve pruebas \texttt{pytest} en
\texttt{tests/test\_cltv.py} y \texttt{tests/test\_dwh\_integrity.py}
(integridad referencial, unicidad de \emph{business keys}, conservación de
ingresos \texttt{staging} $\to$ \texttt{dwh}, y consistencia algebraica de
la fórmula de CLTV).
